\documentclass[]{article}
\usepackage[margin=1in]{geometry}
\usepackage{amsmath, amssymb}
\usepackage{xcolor}
\usepackage[most]{tcolorbox}
\usepackage{enumitem}
\usepackage{fancyhdr}

\geometry{margin=1in}
\pagestyle{fancy}
\fancyhf{}
\rhead{AMC12A 2016 Problem 18 Analysis}
\lhead{\today}

% Color definitions
\definecolor{problemconfig}{RGB}{230, 242, 255}
\definecolor{strategic}{RGB}{255, 230, 242}
\definecolor{tactical}{RGB}{242, 255, 230}
\definecolor{techniques}{RGB}{255, 242, 230}
\definecolor{verification}{RGB}{255, 230, 230}
\definecolor{solution}{RGB}{230, 255, 255}
\definecolor{competition}{RGB}{242, 242, 255}
\definecolor{gold}{RGB}{218, 165, 32}

% Box styles
\newtcolorbox{problemconfigbox}{
    colback=problemconfig,
    colframe=blue,
    title=Problem Configuration Analysis,
    fonttitle=\bfseries,
    arc=3pt,
    breakable,
    boxrule=1pt
}

\newtcolorbox{strategicbox}{
    colback=strategic,
    colframe=purple,
    title=Strategic Framework Application,
    fonttitle=\bfseries,
    arc=3pt,
    breakable,
    boxrule=1pt
}

\newtcolorbox{tacticalbox}{
    colback=tactical,
    colframe=green,
    title=Tactical Methodology Selection,
    fonttitle=\bfseries,
    arc=3pt,
    breakable,
    boxrule=1pt
}

\newtcolorbox{techniquesbox}{
    colback=techniques,
    colframe=orange,
    title=Conversion Techniques Application,
    fonttitle=\bfseries,
    arc=3pt,
    breakable,
    boxrule=1pt
}

\newtcolorbox{verificationbox}{
    colback=verification,
    colframe=red,
    title=Verification Strategy Design,
    fonttitle=\bfseries,
    arc=3pt,
    breakable,
    boxrule=1pt
}

\newtcolorbox{solutionbox}{
    colback=solution,
    colframe=cyan,
    title=Solution Roadmap,
    fonttitle=\bfseries,
    arc=3pt,
    boxrule=1pt,
    breakable
}

\newtcolorbox{competitionbox}{
    colback=competition,
    colframe=purple,
    title=Competition Strategy Integration,
    fonttitle=\bfseries,
    arc=3pt,
    breakable,
    boxrule=1pt
}

\newtcolorbox{keyinsightbox}{
    colback=yellow!20,
    colframe=gold,
    title=Key Insight,
    fonttitle=\bfseries,
    arc=3pt,
    boxrule=2pt,
    leftrule=5pt,
    breakable
}

\newtcolorbox{warningbox}{
    colback=red!10,
    colframe=red!50,
    title=Warning: Common Pitfall,
    fonttitle=\bfseries,
    arc=3pt,
    boxrule=2pt,
    leftrule=5pt,
    breakable
}

\newtcolorbox{tipbox}{
    colback=green!10,
    colframe=green!50,
    title=Pro Tip,
    fonttitle=\bfseries,
    arc=3pt,
    boxrule=2pt,
    leftrule=5pt,
    breakable
}

\begin{document}

\title{AMC12A 2016 Problem 18: Strategic Analysis}
\author{Comprehensive Framework Application}
\date{\today}
\maketitle

\section*{Problem Statement}

For some positive integer $n$, the number $110n^3$ has $110$ positive integer divisors, including $1$ and the number $110n^3$. How many positive integer divisors does the number $81n^4$ have?

\textbf{Answer Choices:} $\textbf{(A) }110\qquad\textbf{(B) }191\qquad\textbf{(C) }261\qquad\textbf{(D) }325\qquad\textbf{(E) }425$

\begin{problemconfigbox}
\subsection*{A. Problem Classification}
\begin{itemize}[leftmargin=*]
    \item \textbf{Problem Type}: To Find - the number of divisors of $81n^4$
    \item \textbf{Difficulty Band}: Late-band (Problem 18 of 25), medium-high difficulty
    \item \textbf{Competition Context}: AMC12 (75 minutes, 25 problems) - approximately 3 minutes per problem
    \item \textbf{Mathematical Domain}: Number Theory (prime factorization, divisor counting)
\end{itemize}

\subsection*{B. Problem Structure Identification}
\begin{itemize}[leftmargin=*]
    \item \textbf{Given Information}: 
    \begin{itemize}
        \item $110n^3$ has exactly $110$ positive integer divisors
        \item $110 = 2 \cdot 5 \cdot 11$
        \item $81 = 3^4$
        \item $n$ is a positive integer
    \end{itemize}
    \item \textbf{Constraints}: 
    \begin{itemize}
        \item $n$ must be chosen such that divisor count condition is satisfied
        \item $n$ likely has prime factorization involving small primes
    \end{itemize}
    \item \textbf{Objective}: Find the number of divisors of $81n^4$
    \item \textbf{Hidden Assumptions}: 
    \begin{itemize}
        \item The divisor formula $(e_1+1)(e_2+1)\cdots(e_k+1)$ applies for $p_1^{e_1}p_2^{e_2}\cdots p_k^{e_k}$
        \item $n$ must be uniquely determined (or at least the form of $n$ is)
    \end{itemize}
    \item \textbf{Answer Choice Leverage}: All answers are distinct; $325 = 5 \times 5 \times 13$ suggests a three-factor product
\end{itemize}

\subsection*{C. Strategic Orientation}
\begin{itemize}[leftmargin=*]
    \item \textbf{Hypothesis}: $110n^3 = 2 \cdot 5 \cdot 11 \cdot n^3$ has $110$ divisors
    \item \textbf{Conclusion}: Determine divisor count of $81n^4 = 3^4 \cdot n^4$
    \item \textbf{Notation}: Let $n = 2^a \cdot 5^b \cdot 11^c \cdot m$ where $m$ is coprime to $2, 5, 11$
    \item \textbf{Intuitive Approach}: Use divisor counting formula; factor $110$ to understand exponent structure
    \item \textbf{Cross-Domain Potential}: Pure number theory; no geometric or algebraic alternative
\end{itemize}
\end{problemconfigbox}

\begin{strategicbox}
\subsection*{A. Psychological Strategies}
\begin{itemize}[leftmargin=*]
    \item \textbf{Mental Toughness}: This is a late-band problem requiring careful factorization and case analysis. Persist through the algebra.
    \item \textbf{Creativity Cultivation}: Recognize that $110 = 2 \times 5 \times 11$ gives structure to the problem
    \item \textbf{Iterative Refinement}: Start with simple cases ($n=1$), then determine what additional factors $n$ needs
\end{itemize}

\subsection*{B. Investigation Strategies}
\begin{itemize}[leftmargin=*]
    \item \textbf{Startup Orientation}: Identify this as a divisor-counting problem using Fundamental Theorem of Arithmetic
    \item \textbf{Get Your Hands Dirty}: Try $n=1$: $110 \cdot 1^3 = 2 \cdot 5 \cdot 11$ has $(1+1)(1+1)(1+1) = 8$ divisors. Need more!
    \item \textbf{Penultimate Step}: Once we find $n$, applying the divisor formula to $81n^4$ is straightforward
    \item \textbf{Wishful Thinking}: Wish that $n$ has a simple form like $2^a$ or $2^a \cdot 5^b$
    \item \textbf{Make It Easier}: Factor $110 = 2 \times 5 \times 11$ to understand the constraint
    \item \textbf{Conjecture via Small Cases}: Build up $n$ systematically to reach $110$ divisors
    \item \textbf{Visualization}: Not applicable (pure number theory)
\end{itemize}

\subsection*{C. Late-Band Strategic Playbook}
\begin{itemize}[leftmargin=*]
    \item \textbf{Answer-Choice Leverage}: $325 = 5^2 \times 13$ suggests final form $3^4 \cdot p_1^{e_1} \cdot p_2^{e_2}$ with $(5)(5)(13)$ structure
    \item \textbf{Make It Easier}: Start with $n=1$, then systematically add prime factors to reach target
    \item \textbf{Penultimate Step}: Key is determining $n$; rest is mechanical
    \item \textbf{Conjecture via Small Cases}: Test $n = 2^3, 2^3 \cdot 5, 2^3 \cdot 5 \cdot 11$ etc.
\end{itemize}

\subsection*{D. Argument Construction Strategy}
\begin{itemize}[leftmargin=*]
    \item \textbf{Direct Proof}: Use divisor formula directly once $n$ is determined
    \item \textbf{Proof by Contradiction}: Not needed
    \item \textbf{Constructive Proof}: Build $n$ step by step to satisfy the constraint
\end{itemize}
\end{strategicbox}

\begin{tacticalbox}
\subsection*{A. Core Mathematical Tactics}
\begin{itemize}[leftmargin=*]
    \item \textbf{Primary Tactic}: Fundamental Theorem of Arithmetic + Divisor Counting Formula
    \item \textbf{Symmetry}: Not applicable
    \item \textbf{Invariants}: The divisor formula structure $(e_1+1)(e_2+1)\cdots$
    \item \textbf{Extremal Principle}: Not directly applicable
    \item \textbf{Pigeonhole Principle}: Not applicable
    \item \textbf{Induction}: Not needed
    \item \textbf{Counting Techniques}: Divisor counting via prime factorization
\end{itemize}

\subsection*{B. Crossover Tactics}
\begin{itemize}[leftmargin=*]
    \item \textbf{Graph Theory}: Not applicable
    \item \textbf{Complex Numbers}: Not applicable
    \item \textbf{Generating Functions}: Not applicable
\end{itemize}
\end{tacticalbox}

\begin{techniquesbox}
\subsection*{A. High-Probability Techniques (Recognition Index)}

\textbf{Applicable Techniques:}
\begin{itemize}[leftmargin=*]
    \item \textbf{TH1: Prime Factorization \& FTA} - \textcolor{red}{CRITICAL}
    \begin{itemize}
        \item Recognition: Problem explicitly involves divisor counting
        \item Application: Write $110n^3 = 2^{1+3a} \cdot 5^{1+3b} \cdot 11^{1+3c} \cdot \text{other primes}$
        \item Why it works: Divisor formula requires complete prime factorization
    \end{itemize}
    
    \item \textbf{TH2: Divisor Sum \& Count Formulas} - \textcolor{red}{CRITICAL}
    \begin{itemize}
        \item Recognition: Problem asks for "number of divisors"
        \item Application: For $p_1^{e_1} \cdots p_k^{e_k}$, divisor count is $(e_1+1) \cdots (e_k+1)$
        \item Why it works: This is the standard divisor counting formula
    \end{itemize}
    
    \item \textbf{TH8: Modular Arithmetic} - Minor role
    \begin{itemize}
        \item Recognition: Working with factorization constraints
        \item Application: Understanding that exponents modulo certain values matter
    \end{itemize}
\end{itemize}

\subsection*{B. Technique Integration}
\begin{itemize}[leftmargin=*]
    \item \textbf{Primary Technique}: FTA + Divisor Formula (TH1 + TH2)
    \item \textbf{Supporting Techniques}: Case analysis, systematic enumeration
    \item \textbf{Technique Order}:
    \begin{enumerate}
        \item Factor $110 = 2 \times 5 \times 11$
        \item Determine that $110n^3$ has at least 3 prime factors
        \item Use divisor formula: $(e_1+1)(e_2+1)(e_3+1) \cdots = 110$
        \item Find valid exponent combinations
        \item Determine $n$'s form
        \item Apply divisor formula to $81n^4 = 3^4 \cdot n^4$
    \end{enumerate}
\end{itemize}

\subsection*{C. Rationale for Selection}
This is a pure number theory problem requiring fundamental divisor counting techniques. The key insight is recognizing that $110 = 2 \times 5 \times 11$ constrains the prime factorization structure of $n$.
\end{techniquesbox}

\begin{verificationbox}
\subsection*{A. Verification Level Selection}

\textbf{Recommended Level}: Level 4 (Standard) to Level 5 (Thorough)

\textbf{Decision Tree Path}:
\begin{itemize}
    \item Problem difficulty: Late-band (P18) $\rightarrow$ High stakes
    \item Domain: Number theory $\rightarrow$ Calculation errors possible
    \item Confidence after solve: Medium-High $\rightarrow$ Verify thoroughly
    \item Time remaining: If adequate, use Level 5
\end{itemize}

\subsection*{B. Verification Methods}
\begin{itemize}[leftmargin=*]
    \item \textbf{Primary Check}: Verify that $110n^3$ with one example form of $n$ actually has $110$ divisors
    \begin{itemize}
        \item If $n = 2^3 \cdot 5$, then $110n^3 = (2 \cdot 5 \cdot 11)(2^3 \cdot 5)^3 = 2 \cdot 5 \cdot 11 \cdot 2^9 \cdot 5^3 = 2^{10} \cdot 5^4 \cdot 11^1$
        \item Divisor count: $(10+1)(4+1)(1+1) = 11 \times 5 \times 2 = 110$ $\checkmark$
    \end{itemize}
    
    \item \textbf{Secondary Check}: Verify divisor count of $81n^4$ (CRITICAL - this is what the problem asks!)
    \begin{itemize}
        \item $81n^4 = 3^4 \cdot (2^3 \cdot 5)^4 = 3^4 \cdot 2^{12} \cdot 5^4$
        \item Divisor count: $(4+1)(12+1)(4+1) = 5 \times 13 \times 5 = 325$ $\checkmark$
    \end{itemize}
    
    \item \textbf{Answer Choice Verification}: $325 = \boxed{\textbf{(D)}}$ matches
    
    \item \textbf{Structure Check}: Verify $n$ has form $p_1 \cdot p_2^3$
    \begin{itemize}
        \item All 6 permutations of exponents $\{1,4,10\}$ among primes $\{2,5,11\}$ give $n = p_1 \cdot p_2^3$ form
        \item This means $n^4 = p_1^4 \cdot p_2^{12}$, which combined with $3^4$ gives $(5)(5)(13) = 325$ divisors
        \item The specific permutation doesn't matter—all lead to same answer!
    \end{itemize}
\end{itemize}

\subsection*{C. Confidence Calibration}
\begin{itemize}[leftmargin=*]
    \item \textbf{Confidence Level}: 95\%
    \item \textbf{Rationale}: Method is standard, calculations verified, answer matches expected form
    \item \textbf{Flag for Review}: No, unless time permits final recheck
\end{itemize}
\end{verificationbox}

\begin{solutionbox}
\subsection*{Solution Roadmap}

\textbf{Total Estimated Time}: 4-5 minutes

\subsubsection*{Step 1: Analyze the Given Information (30 seconds)}
\begin{itemize}
    \item Given: $110n^3$ has $110$ divisors
    \item Factor: $110 = 2 \cdot 5 \cdot 11$
    \item Note: $110n^3 = 2 \cdot 5 \cdot 11 \cdot n^3$
\end{itemize}

\subsubsection*{Step 2: Apply Divisor Formula (1 minute)}
\begin{itemize}
    \item For $n=1$: $110 = 2 \cdot 5 \cdot 11$ has $(1+1)(1+1)(1+1) = 8$ divisors
    \item Need $110$ divisors, so $110 = 2 \times 5 \times 11$ must be factored
    \item Since $110n^3$ has at least three distinct primes $(2, 5, 11)$, and divisor formula gives products, we need $(a+1)(b+1)(c+1) = 110$
\end{itemize}

\subsubsection*{Step 3: Determine Exponent Structure (1.5 minutes)}
\begin{itemize}
    \item $110 = 2 \times 5 \times 11$ (only way to write as product of 3 factors $\geq 2$)
    \item So we need exponents such that $(e_1+1)(e_2+1)(e_3+1) = 2 \times 5 \times 11$
    \item The exponents are $\{1, 4, 10\}$ in some order
    \item Since $110n^3 = 2^{1+3a} \cdot 5^{1+3b} \cdot 11^{1+3c}$, where $a,b,c \geq 0$
    \item Testing: If exponent of $2$ is $1$, then $1+3a = 1 \Rightarrow a = 0$
    \item If exponent of $2$ is $4$, then $1+3a = 4 \Rightarrow a = 1$ \checkmark
    \item If exponent of $2$ is $10$, then $1+3a = 10 \Rightarrow a = 3$ \checkmark
    \item Similarly for $5$ and $11$
\end{itemize}

\subsubsection*{Step 4: Deduce Form of $n$ (1 minute)}
\begin{itemize}
    \item Key insight from Solution 2: $110n^3$ has exactly 3 prime factors (2, 5, 11) with divisor count $110 = 2 \times 5 \times 11$
    \item Since $(e_1+1)(e_2+1)(e_3+1) = 110$ and none can be 1, the exponents are $\{1, 4, 10\}$ in some order
    \item Consider $110n^3 = 2^{a} \cdot 5^{b} \cdot 11^{c}$ where exponents are from $\{1, 4, 10\}$
    \item Testing: If $a=10, b=4, c=1$, then $110n^3 = 2^{10} \cdot 5^4 \cdot 11^1$
    \item Since $110 = 2^1 \cdot 5^1 \cdot 11^1$, we need $n^3$ to contribute: $2^9, 5^3, 11^0$
    \item Therefore: $n^3 = 2^9 \cdot 5^3 \Rightarrow n = 2^3 \cdot 5$
    \item Verification: $110n^3 = (2 \cdot 5 \cdot 11)(2^3 \cdot 5)^3 = 2 \cdot 5 \cdot 11 \cdot 2^9 \cdot 5^3 = 2^{10} \cdot 5^4 \cdot 11^1$ $\checkmark$
\end{itemize}

\textbf{Key observation}: There are actually 6 possible forms for $n$ (all permutations of assigning exponents $\{1,4,10\}$ to primes $\{2,5,11\}$), but they all lead to $n = p_1 \cdot p_2^3$ structure where $p_1, p_2 \in \{2, 5, 11\}$

\subsubsection*{Step 5: Calculate Divisors of $81n^4$ (1 minute)}
\begin{itemize}
    \item $81n^4 = 3^4 \cdot n^4 = 3^4 \cdot (2^3 \cdot 5)^4 = 3^4 \cdot 2^{12} \cdot 5^4$
    \item Divisor count: $(4+1)(12+1)(4+1) = 5 \times 13 \times 5 = 325$
    \item Answer: $\boxed{\textbf{(D) } 325}$
\end{itemize}

\subsubsection*{Checkpoints \& Alternatives}
\begin{itemize}
    \item \textbf{Checkpoint 1}: Verified $110 = 2 \times 5 \times 11$ is correct factorization
    \item \textbf{Checkpoint 2}: Confirmed exponent structure matches $n$ form
    \item \textbf{Alternative}: Could verify by computing divisors of $81 \cdot 40^4$ directly (tedious)
    \item \textbf{Pivot Indicator}: If initial $n$ guess doesn't yield $110$ divisors, reassess exponent distribution
\end{itemize}
\end{solutionbox}

\begin{competitionbox}
\subsection*{Competition Strategy Integration}

\subsubsection*{A. Time Management}
\begin{itemize}[leftmargin=*]
    \item \textbf{Target Time}: 4-5 minutes (Problem 18, late-band)
    \item \textbf{Time Allocation}:
    \begin{itemize}
        \item Reading \& Setup: 30 seconds
        \item Work \& Calculation: 3 minutes
        \item Verification: 1-1.5 minutes
    \end{itemize}
    \item \textbf{Actual vs. Expected}: Should complete within target; if exceeding 6 minutes, flag and move on
\end{itemize}

\subsubsection*{B. Problem Prioritization}
\begin{itemize}[leftmargin=*]
    \item \textbf{Accessibility}: Moderate-high. Requires solid number theory knowledge
    \item \textbf{Point Value}: 6 points (same as all AMC12 problems)
    \item \textbf{Priority Level}: Medium-high. Late-band problems are challenging but doable with proper technique
\end{itemize}

\subsubsection*{C. Risk Management}
\begin{itemize}[leftmargin=*]
    \item \textbf{Confidence Threshold}: 85\%+ to submit
    \item \textbf{Error Recovery}: If calculation error detected, restart from prime factorization
    \item \textbf{Abandonment Criteria}:
    \begin{itemize}
        \item If unable to factor $110$ correctly after 2 minutes
        \item If confusion persists about exponent structure after 3 minutes
        \item Hard abandon if approaching 7+ minutes with no clear path
    \end{itemize}
\end{itemize}

\subsubsection*{D. Answer Format}
\begin{itemize}[leftmargin=*]
    \item Multiple choice: Bubble (D) carefully
    \item Quick sanity check: $325$ is reasonable for a late-band answer
    \item Elimination strategy: Answers (A) $110$ and (B) $191$ seem too low given $81n^4$ likely has more divisors than $110n^3$
\end{itemize}
\end{competitionbox}

\begin{keyinsightbox}
\textbf{Key Insight}: The constraint that $110n^3$ has exactly $110$ divisors, combined with the factorization $110 = 2 \times 5 \times 11$, determines that $110n^3$ must have exactly 3 distinct prime factors with exponents that yield $(e_1+1)(e_2+1)(e_3+1) = 110$. Since 110 cannot be written as a product of three factors $\geq 2$ in any other way, the exponents must be $\{1, 4, 10\}$ in some order. This forces $n = p_1 \cdot p_2^3$ where $p_1, p_2 \in \{2,5,11\}$. Crucially, ANY such permutation leads to the same answer for $81n^4$, since the structure $3^4 \cdot p_1^4 \cdot p_2^{12}$ always gives $(5)(5)(13) = 325$ divisors!
\end{keyinsightbox}

\begin{warningbox}
\textbf{Warning}: Do not confuse the exponent structure! When $110n^3 = 2^{10} \cdot 5^4 \cdot 11$, remember that the base factorization is $2 \cdot 5 \cdot 11 \cdot n^3$, so you must carefully extract what powers of each prime are in $n$. Also, ensure you raise $n$ to the correct power when computing $81n^4$.
\end{warningbox}

\begin{tipbox}
\textbf{Pro Tip}: When solving divisor problems, always factor the given numbers first ($110 = 2 \times 5 \times 11$ and $81 = 3^4$). This immediately reveals the prime structure and guides your analysis. For late-band AMC problems, recognize that answer choices often hint at the structure—here, $325 = 5^2 \times 13$ suggests a clean three-factor divisor formula.
\end{tipbox}

\section*{Final Answer}
\[\boxed{\textbf{(D) } 325}\]

\section*{Confidence Assessment}
\begin{itemize}[leftmargin=*]
    \item \textbf{Confidence Level}: 95\%
    \item \textbf{Justification}: Standard number theory technique applied correctly; calculations verified; answer matches expected structure
    \item \textbf{Verification Applied}: Level 5 (Thorough) - checked divisor counts for both $110n^3$ and $81n^4$
    \item \textbf{Flag for Review}: No
    \item \textbf{Time Spent}: Estimated 4-5 minutes (within target for P18)
\end{itemize}

\end{document}